\subsection*{MISC Library}

The following
names are provided by the MISC library:
\begin{itemize}
  \item \verb#int#\texttt{(x)}, \verb#float#\texttt{(x)},
    \verb#bool#\texttt{(x)}, \verb#complex#\texttt{(x)}, \verb#str#\texttt{(x)}: \textit{primitive},
    converts the value \texttt{x} to an integer, a float, a boolean, a complex number, or a string
    representation of \texttt{x} respectively, and returns the
    respective value. For \verb#int#\texttt{(x)}, if \texttt{x} is a string,
    \texttt{x} is parsed
    as an integer in base 10.
  \item \verb#int#\texttt{(s, i)}: \textit{primitive},
    interprets the \emph{string} \texttt{s} as an integer, using the
    positive integer \texttt{i} as radix, and returns the respective value.
  \item \verb#complex#\texttt{(r, i)}: \textit{primitive},
    takes in two numbers \texttt{r} and \texttt{i}, and returns the complex number 
    with real part \texttt{r} and imaginary part \texttt{i}.
  \item \verb#complex#\texttt{()}: \textit{primitive},
    returns the complex number with real part 0 and imaginary part 0.
  \item \verb#repr#\texttt{(s)}: \textit{primitive},
    returns a string representation of \texttt{s} that is unambiguous and suitable for
    debugging; for many values, this string can be used to reconstruct \texttt{s}. 
  \item
    \href{https://sourceacademy.org/sicpjs/4.1.2\#p2}{\lstinline{is_boolean(x)}},
    \href{https://sourceacademy.org/sicpjs/2.3.2\#p5}{\lstinline{is_int(x)}},
    \href{https://sourceacademy.org/sicpjs/2.3.2\#p5}{\lstinline{is_float(x)}},
    \href{https://sourceacademy.org/sicpjs/2.3.2\#p7}{\lstinline{is_string(x)}},
    \href{https://sourceacademy.org/sicpjs/4.1.2\#p2}{\lstinline{is_none(x)}},
    \verb#is_function#\texttt{(x)},
    \verb#is_complex#\texttt{(x)}: \textit{primitive}, returns
    \texttt{True} if the type of \texttt{x} matches the function name
    and \texttt{False} if it does not.
  \item \texttt{real(x)}, \texttt{imag(x)}: \textit{primitive}, return the 
    real and imaginary parts of a complex number \texttt{x} respectively.
  \item \texttt{input(s)}: \textit{primitive}, pops up a window that
    displays the \emph{string} \texttt{s}, provides
    an input line for the user to enter a text, a ``Cancel'' button
    and an ``OK'' button. The call of \texttt{input}
    suspends execution of the program until one of the two buttons is
    pressed. If
    the ``OK'' button is pressed, \texttt{input} returns the entered
    text as a string.
    If the ``Cancel'' button is pressed, \texttt{input} returns a
    non-string value.
  \item
    \verb#min#\texttt{(x, y, *args)} and \verb#max#\texttt{(x, y, *args)}: \textit{primitive}, 
    takes in two or more arguments and returns the smallest and largest of the arguments respectively. 
    If multiple arguments are equally small or large, the first one encountered in the argument list is returned. 
    The arguments must be mutually comparable (i.e., they must support the < and > operators).
  \item
    \verb#abs#\texttt{(x)}: \textit{primitive}, returns the absolute value of \texttt{x}.
    For an \verb#int# input, it returns the non-negative integer equivalent. 
    For a \verb#float# input, it returns the positive floating-point number representing its magnitude.
    For a \verb#complex# input, it returns the modulus of the complex number, which is a non-negative float.
  \item
    \verb#print#\texttt{(*x)}: \textit{primitive}, takes in any number of parameters, converts them to their string
    representations using \verb#str# and writes them to the console, followed by a newline character.
  \item
    \verb#error#\texttt{(*x)}: \textit{primitive}, takes in any number of parameters, converts them to their string
    representations using \verb#str# and writes them to the console, followed by a newline character.
    The evaluation of any call of \texttt{error} aborts the running program immediately.
  \item
    \verb#arity#\texttt{(f)}: \textit{primitive}, 
    takes in a function \texttt{f} and returns the minimum number of parameters that \texttt{f} can accept.
  \item
    \verb#round#\texttt{(x)}\footnote{
      In \href{https://docs.python.org/3.13/library/functions.html\#round}{Python 3.13},
      when the value to be rounded is exactly halfway between two possible rounded values, it rounds to the nearest even number. For example, \texttt{round(2.5)} would return 2, while \texttt{round(3.5)} would return 4.
    }: \textit{primitive}, 
    takes in an int or float \texttt{x} and returns the nearest integer to \texttt{x}.
  \item
    \verb#round#\texttt{(x, n)}: \textit{primitive},
    takes in an int or float \texttt{x} and an int \texttt{n}, 
    and returns the value of \texttt{x} rounded to \texttt{n} digits of precision after the decimal point.
    If the second argument \texttt{n} is None, it is treated as \verb#round#\texttt{(x)}.
  \item
    \verb#len#\texttt{(s)}: \textit{primitive},
    takes in a container for a finite number of values (e.g., a string), and returns the number of elements in the container.
  \item
    \verb#time_time#\texttt{()}: \textit{primitive},
    returns number of milliseconds elapsed since
    January 1, 1970 00:00:00 UTC
  \item
    \verb#random_random#\texttt{()}: \textit{primitive},
    returns a random float uniformly in the range $[0.0, 1.0)$.
  

\end{itemize}
All library functions can be assumed to run
in $O(1)$ time, except \texttt{print}, \texttt{error}, \texttt{str} and \texttt{repr},
which run in $O(n)$ time, where $n$ is
the size (number of components such as pairs) of their arguments.
